\documentclass[10pt]{article}
\usepackage[margin=.8in]{geometry}
\usepackage{amsmath,hyperref}
\title{Guide: certified request cutoffs}
\author{Collatz Lean research project}
\date{September 6, 2026}
\begin{document}
\maketitle
\section*{What is now checked}
The earlier optimizer proposed a seed cutoff from the transfer parameters available in an induction. The new Lean checker verifies that no odd-run request below that cutoff needs an unavailable parameter.

For transfer parameters below 61, the certified seed cutoff is 415, improving the sufficient dyadic cutoff 128. For parameters below $2^{64}-1$, the certified cutoff is 30786325577727. These are conditional convergence thresholds: all smaller transfer instances must still be available.
\section*{Using the certificate}
Choose U, a proposed seed bound N, and L such that $N<2^L$.
Prove \texttt{RequestCutoff.Check U N L}. The checker considers each run length below L and the two odd residue classes modulo four. It checks only their largest eligible odd parts.

Use \texttt{RequestCutoff.reachesOne} to obtain convergence below N from all transfer instances below U. To close a transfer target that descends into this interval, use
\texttt{AffineTransfer.of\_request\_cutoff\_descent}.
The source convergence argument alone does not establish the cutoff's smaller-transfer premises.
\section*{Source and verification}
\raggedright
The formal source is \texttt{lean/Collatz/RequestCutoff.lean}.
The declarations \texttt{check\_61} and \texttt{check\_mersenne64} are evaluated by ordinary \texttt{decide}, with no extra computational axiom.
The declarations \texttt{maximal\_61} and \texttt{maximal\_mersenne64} exhibit the first unavailable requests and exclude every larger cutoff in the same request system.

From \texttt{lean/}, run:
\begin{verbatim}
lake build Collatz.RequestCutoff
\end{verbatim}
From the repository root, with the pinned toolchain on PATH:
\begin{verbatim}
python3 analysis/audit_theorems.py --build
python3 analysis/transfer_requests.py
python3 -m unittest discover -s analysis \
  -p 'test_transfer_requests.py'
\end{verbatim}
The Python program proposes thresholds. Only the two selected outputs described here are certified in this module; its full enumeration and optimizer are not formalized as Lean algorithms.
\section*{The remaining mathematical task}
A cutoff certificate supplies an induction premise, not a descent witness. We still need a general argument that targets descend into their certified intervals, or a covering combination of this rule and other valid transfer rules. Maximality of these cutoffs does not establish nonconvergence above them. The Collatz conjecture remains unproved.
\end{document}
